\chapter{Combinatoria}

Aunque la enumeración o conteo se percibe como un concepto básico de la aritmética, su importancia y complejidad suelen subestimarse. Es esencial comprender los principios fundamentales en lugar de confiar únicamente en fórmulas.

\section{Principio de la multiplicación}

Si una actividad se puede realizar en $t$ pasos sucesivos, y el paso $i$ se puede hacer de $n_i$ maneras, entonces el número total de actividades es $n_1 \cdot n_2 \cdots n_t$.

\begin{ejemplo}
En una picantería se ofrece: 2 entradas, 3 platos de fondo y 4 bebidas. ¿Cuántas formas hay de seleccionar un plato de fondo y una bebida?
\[
\text{Respuesta} = 3 \cdot 4 = 12.
\]
¿Y una entrada, un plato y una bebida?
\[
\text{Respuesta} = 2 \cdot 3 \cdot 4 = 24.
\]
\end{ejemplo}

\begin{ejemplo}
¿Cuántas cadenas de longitud 4 se pueden formar con las letras ABCDE sin repetición?
\begin{enumerate}[label=\alph*)]
  \item Total: $5 \cdot 4 \cdot 3 \cdot 2 = 120$.
  \item Comienzan con B: $1 \cdot 4 \cdot 3 \cdot 2 = 24$.
  \item No comienzan con B: $120 - 24 = 96$.
\end{enumerate}
\end{ejemplo}

\begin{ejemplo}
¿Cuántos subconjuntos tiene $A = \{a_1, \ldots, a_n\}$?
\[
\underbrace{2 \cdot 2 \cdots 2}_{n \text{ factores}} = 2^n.
\]
\end{ejemplo}

\section{Principio de la Suma}

Si $X_1, \ldots, X_t$ son conjuntos ajenos por pares con $n_i$ elementos cada uno, el número de elementos en la unión es $n_1 + \cdots + n_t$.

\begin{ejemplo}
¿Cuántas cadenas de ocho bits comienzan con 101 o con 111?
\[
2^5 + 2^5 = 64.
\]
\end{ejemplo}

\begin{definicion}
Una \textbf{permutación} de $n$ elementos distintos es un ordenamiento de los $n$ elementos.
\end{definicion}

\begin{teorema}
Existen $n!$ permutaciones de $n$ elementos.
\end{teorema}
\begin{proof}
Por el principio de la multiplicación:
$n(n-1)(n-2)\cdots 3\cdot 2\cdot 1 = n!$.
\end{proof}

\begin{definicion}
Una \textbf{permutación $r$} de $n$ elementos es un ordenamiento de $r$ de ellos. Se denota $P(n,r)$.
\end{definicion}

\begin{teorema}
\[
P(n,r) = \frac{n!}{(n-r)!}, \quad r \leq n.
\]
\end{teorema}

\begin{ejemplo}
¿De cuántas maneras se puede elegir presidente, vicepresidente, secretario y tesorero entre 10 delegados?
\[
P(10,4) = 10\cdot 9\cdot 8\cdot 7 = 5040.
\]
\end{ejemplo}

\begin{definicion}
Una \textbf{combinación $r$} de $X$ es una selección no ordenada de $r$ elementos de $X$. Se denota $C(n,r)$ o $\binom{n}{r}$.
\end{definicion}

\begin{teorema}
\[
C(n,r) = \frac{n!}{(n-r)!\,r!}.
\]
\end{teorema}

\begin{ejemplo}
¿De cuántas maneras elegir 3 representantes entre 5 delegados?
\[
C(5,3) = \frac{5!}{2!\,3!} = 10.
\]
\end{ejemplo}

\begin{ejemplo}
Comité de 2 mujeres y 3 hombres de 5 mujeres y 6 hombres:
\[
C(5,2)\,C(6,3) = 10 \cdot 20 = 200.
\]
\end{ejemplo}

\begin{ejemplo}
¿Cuántas cadenas de 8 bits contienen exactamente 4 unos?
\[
C(8,4) = 70.
\]
\end{ejemplo}

\begin{ejemplo}
Baraja de 52 cartas:
\begin{enumerate}[label=\alph*)]
  \item Manos de póquer: $C(52,5) = 2\,598\,960$.
  \item Todas del mismo palo: $4 \cdot C(13,5) = 5148$.
  \item Tres de una denominación y dos de otra: $13 \cdot 12 \cdot C(4,3) \cdot C(4,2) = 3744$.
\end{enumerate}
\end{ejemplo}